For a local observable $F_{N,p}(y)$, let $a_N^eF_{N,p}$ be its dimensionless representative. We require it to be continuous on its fixed stencil, regular on healthy configurations in the sense of Definition~\ref{def:D05a}, with derivative and supremum bounds of subpower growth. The supremum bound is on all configurations. These two bounds have different uses: Taylor extraction uses the first, and the forced complement uses the second.

Local matching makes no assumption on the size of the physical coefficient it produces. It bounds its inverse terminal coordinate, which is the quantity in~\eqref{eq:F-subpower}.
\begin{lemma}[Local comparison]\label{lem:E-local-comparison}
For every such degree-$e$ family $F$, there are admissible coefficients $b_{N,p}\in V_e$ with $O^{(e)}(b)-F$ small. For the auxiliary section of Lemma~\ref{lem:A-R4}, there are invertible coefficient changes $P_{N,p}$, with $\|P_{N,p}^{-1}\|_1\le2$ eventually uniformly in~$p$, and maps
\begin{equation*}
 B_{d,N,p}=\mathsf Z_{N,p,d}P_{N,p}\factorscale_{t_N}\mathsf R_d
\end{equation*}
from polynomial labels to $V_d$, such that $O^{(d)}(B_{d,N,p}\pi)-a_N^{-d}V_d^{\rm aux}[\mathsf R_d\pi]$ is small for equation labels with admissible $B_{d,N,p}\pi$.
\end{lemma}
In particular, for $\pi_{N,p}\in\mathcal I_d$ with admissible $B_{d,N,p}\pi_{N,p}$, the comparison is
\begin{equation}\label{eq:E-equation-comparison}
 O^{(d)}_{N,p}(B_{d,N,p}\pi_{N,p},y)
 -a_N^{-d}V_d^{\rm aux}[\mathsf R_d\pi_{N,p}](y)
 \quad\text{is small}.
\end{equation}
We prove both assertions using the same backward coefficient correction.
\begin{proof}
Use Definitions~\ref{def:D05a}, \ref{def:D07}, \ref{def:D08}, Lemmas~\ref{lem:E-marked-step}, \ref{lem:E-replication} and the local sections of Lemma~\ref{lem:A1}. Write $h_N=a_N^eF_{N,p}(0)$. Enlarge its stencil by the section stencil. The radial Taylor readout through degree~$e$ gives a coefficient $v_N$ and a remainder. The readout is bounded by a fixed multiple of the derivative bound: it is a linear map on a finite jet space. Subtract the section with that coefficient on the healthy part. Its regular remainder has zero value and zero radial jet through degree~$e$. Assign the complementary observable to the forced supports meeting the enlarged stencil. This gives an exact one-mark representation of $\density\rho_N h_N$, with
\begin{equation}\label{eq:E-local-start}
 |v_N|_1+r_N\le E_{N,p},\qquad
 \text{forced majorant}\le E'_{N,p},
\end{equation}
where $E,E'$ have subpower growth. The support weight of this fixed enlarged stencil is bounded by $e^{81\sigma_{J_N}}$ eventually; its $B_N$ covering number is at most~$81$. Threshold powers, the readout bound and this support weight are included in $E,E'$. The entry conditions make each of them subpower. No division by the forced allowance is made on the uncontrolled complement: it is retained with its original eligibility tests and the supremum bound $E'$.

To compare with the selected graph, carry this representation backwards and at the reference level choose the primary coefficient with the same polynomial coordinate. Subtract the corresponding selected inserted density. Its polynomial coordinate at the reference level is zero. The following calculation carries this matching condition towards the fine scale.

For one step, let $u,v$ be the coarse and fine coordinates, $R,T$ their remainder bounds, $Y=|\mathsf S_Lv|_1$, $K$ the norm bound for the leading map and its inverse, $\varepsilon_L=L^{-3/8}$ and $\vartheta_L=L^{1/8}$. The marked-step inequalities give
\begin{equation}\label{eq:E-match-cone}
 Y\le K\bigl(|u|_1+\varepsilon_L(Y+T)\bigr),\qquad R\le KY+\varepsilon_L T.
\end{equation}
Suppose $a|u|_1\le R$, with $a=32(1+C_{J_i})^2\ge4K^2$. Then $4K|u|_1\le Y+\varepsilon_L T$. Substitution in the first inequality, using $4K\varepsilon_L\le1$, gives $Y\le4K\varepsilon_L T$. Item~\ref{it:E-match} in Lemma~\ref{lem:E1.1} gives
\begin{equation*}
 4K\,32(1+C_{J_{i+1}})^2\varepsilon_L\le1,\qquad
 (4K^2+1)\vartheta_L\varepsilon_L\le1.
\end{equation*}
Consequently $32(1+C_{J_{i+1}})^2|v|_1\le T$ and $\vartheta_L R\le T$. The matching condition starts this induction with $u=0$ at the reference level. Multiplying the second inequality through the interval gives the inverse-scale gain, while the finest-level bound is closed by absorbing the coefficient contribution into its remainder bound. In particular, with subpower $E''$ absorbing $E,E'$,
\begin{equation}\label{eq:E-backward-match}
 |\Delta_N|_1\le\frac{2(1+C_{J_N})E''_{N,p}}{\ell_{N-1}^{1/8}},\qquad
 R_k\le2(1+C_{J_N})E''_{N,p}e^{-(s_N-s_k)/8}.
\end{equation}
Here $\Delta_N$ is the difference from the selected fine polynomial coordinate. The same bounds hold when a local source is first represented at any level $n\le N$, with its polynomial coordinate matched to the selected graph at a lower reference level. If $E''_{n;N,p}$ bounds its starting data after the coefficient contribution is absorbed, then
\begin{equation}\label{eq:E-backward-match-level}
 |\Delta_n|_1\le\frac{2(1+C_{J_n})E''_{n;N,p}}{\ell_{n-1}^{1/8}},\qquad
 R_k\le2(1+C_{J_n})E''_{n;N,p}e^{-(s_n-s_k)/8}
 \quad(r(e)\le k<n).
\end{equation}
Indeed the cone starts at the same zero matched coordinate, and its induction uses only the steps through~$n$. The boundary absorption uses the starting bound there. The product $\prod_{i=k}^{n-1}\ell_i^{-1/8}$ is exactly $e^{-(s_n-s_k)/8}$. Taking $n=N$ recovers~\eqref{eq:E-backward-match}. Thus this factor comes from the matched remainder inequalities, rather than from the physical diameter of the original stencil.

To estimate the replicated residual, for a source of size $E$ started at $N$, use at level~$k$ the radius
\begin{equation*}
 w=\frac{\zeta_{e,k}/4}{a_k^{-e}E\exp\{-(s_N-s_k)/8\}},
\end{equation*}
where $\zeta_{e,k}=\widehat\beta_k^{-(q_e+4)}$ is the source allowance fixed before~$N,p$ in~\eqref{eq:E-attachment-allowance}. Its reciprocal is a fixed power of the reference coupling. At every finer step multiply this radius by the source contraction factor of that step. Its product with the dimensionless singleton size is at most one quarter of the available source allowance; the other three quarters bound the graph remainder, the forced complement and their sum. This follows inductively from the two inequalities in~\eqref{eq:E-atom-estimates}. For the first step the forced complement is bounded by $E'$ in~\eqref{eq:E-local-start}; the large-field suppression absorbs its fixed support and derivative factors before division by the new forced normalization. The radius transport~\eqref{eq:E-source-allowance-transport} and the site-sum calculation in the proof of Lemma~\ref{lem:E-source-attachments} verify the attachment bound at every finer level. Thus every step is admitted. The singleton polynomial coordinate at the reference level is zero, so the representation constructed this way is exactly the difference with the selected primary coefficient, rather than a new approximation to it.

Here $m$ labels the tiling and $k$ the finer stopping level at which its replicated sources are integrated. Replicate the difference at an admissible pair $m\le k$, $B_k^2\le a_m/a_k$, and use the weighted source representation~\eqref{eq:E-source-form}. Coefficient extraction contributes $w^{-1}$ per mark; the complete singleton bounds and the normalized unmarked ratio~\eqref{eq:E-vacuum-ratio} contribute a constant depending on $m,k$. Consequently, for every fixed such pair,
\begin{equation}\label{eq:E-local-error}
 \|O^{(e)}(b)-F\|_{\mathcal C^{\rm tile}_m(0)}
 \le C_{m,k}\mathfrak e_{N,p}a_N^{1/8},\qquad
 \|\mathsf Z_{N,p,e}^{-1}b_{N,p}\|_\infty\le2E'''_{N,p},
\end{equation}
where $\mathfrak e,E'''$ have subpower growth. The factor $e^{s_k/8}$ and the fixed radius $\zeta_{e,k}^{-1}$ belong to $C_{m,k}$. The same argument applies after each signed permutation and conjugation, with translation covariance and the union of their fixed stencils. Absorbing two factors $a_N^{-1/32}$ bounds the first expression by $a_N^{1/16}$ on every fixed pair. The interior-root comparison after Definition~\ref{def:E2-small} gives the common tolerance $a_N^{1/32}$, chosen before all cell requests. This proves smallness and admissibility on full cutoff tails.

For the auxiliary section, perform the construction on all its basis polynomials together, with the original Taylor normalization. Let $u_j$ be the primary bare coefficient matching its $j$th basis vector, and set $Px=\sum_jx_ju_j$, after the finite symmetry averaging. For this fixed section the regular derivative bound $H_{j,d}$ at stage~$j$ is fixed before its ratio, independently of the later coupling. The sharper coefficient estimate in~\eqref{eq:E-backward-match} therefore has numerator bounded by $2(1+C_j)H_{j,d}$ on each basis input. The ratio condition, including the finite basis-recombination factor, gives
\begin{equation*}
 \left|x-Px\right|_1\le\tfrac14|x|_1.
\end{equation*}
Indeed, before summing over the fixed basis the right side is the last-step factor $\ell_{N-1}^{-1/8}$ times the stage coefficient bound; the stage ratio was chosen after that bound. The matrix is therefore invertible and $|P^{-1}x|_1\le2|x|_1$. The construction is common to all later cell requests. Multiplication by $\factorscale_{t_N}\mathsf R_d$ changes a polynomial label into its auxiliary bare insertion coordinate; multiplication by $\mathsf Z$ changes that coordinate to the reference normalization. These are exactly the maps $B$ in the statement. The inverse bound preserves admissibility, and~\eqref{eq:E-local-error} compares the two sections on its transported equation ideal. Finite combinations use the same subpower factors. The argument only replaces eventual bounds on coefficients, so ordered-admissible families obey the corresponding ordered version as well.
\end{proof}

The next positional estimate is used before bounds on differentiated fields have been established. For a grid configuration~$x$ define its minimum coordinate separation on the physical torus by
\begin{equation*}
 g_p(x)=\min_{i\ne j}\min_{a=1,\ldots,4}\min_{z\in L_p\Z}
 |x_i^a-x_j^a+z|,
\end{equation*}
with value $+\infty$ for one slot. This separation supplies the empty slabs for reflection positivity. A coordinate hybrid of $x,y$ chooses each coordinate from either configuration. Assume the insertions and their single-axis reflections are interior at a centered level~$D\le N$, with all their roots at most~$A\ge0$.
\begin{lemma}[Displacement of correlation tests]\label{lem:E-displacement}
For every $r>0$ with $a_D<r$, a one-coordinate grid move $x\mapsto y$ of size at most~$r$ from $g_p(x)\ge4r$ satisfies
\begin{equation*}
 |M_{N,p}(v;x/a_N)-M_{N,p}(v;y/a_N)|
 \le A^n|x-y|/r.
\end{equation*}
If $|x-y|\le r$ and every coordinate hybrid has coordinate separation at least~$4r$, then
\begin{equation*}
 |M_{N,p}(v;x/a_N)-M_{N,p}(v;y/a_N)|
 \le4nA^n|x-y|/r.
\end{equation*}
\end{lemma}
\begin{proof}
Use Definition~\ref{def:D27} and the reflection inequality~\eqref{eq:E-chess}; Lemma~\ref{lem:E2.3} supplies the assumed roots for bounded coefficient lists on full cutoff tails. The separation hypothesis gives every directed coordinate distance at least~$4r$, at any period. No fixed row neighbourhood is used.

Fix an axis and all other coordinates. Reflection positivity and translation invariance give a positive transfer operator, rescaled to $0\le T\le1$. On the periodic axis the expectation is a normalized trace. If $X$ denotes the insertion being moved, its two adjacent empty slabs have lengths $n_1^{\rm sl},n_2^{\rm sl}$ in fine sites. A move by $l\ge0$ changes their factor by
\begin{align}\label{eq:E-trace-displacement}
 T^{n_1^{\rm sl}+l}XT^{n_2^{\rm sl}-l}-T^{n_1^{\rm sl}}XT^{n_2^{\rm sl}}
 &=T^{n_1^{\rm sl}}(T^lX-XT^l)T^{n_2^{\rm sl}-l}\notag\\
 &=-T^{n_1^{\rm sl}}(1-T^l)XT^{n_2^{\rm sl}-l}
   +T^{n_1^{\rm sl}}X(1-T^l)T^{n_2^{\rm sl}-l}.
\end{align}
Both expressions are inserted in the same remaining cyclic product and divided by the same partition trace. The stencil is contained in a centered cell of lattice side $d_{\rm cell}=a_D/a_N$. Since $a_D<r$, the directed separations exceed $2d_{\rm cell}+2r/a_N$, and $la_N\le r$ by assumption. Splitting the two surrounding empty slabs at their reflection planes leaves $n_{\rm sl}+1\ge2r/a_N$ fine sites next to each occurrence of $1-T^l$. The spectral theorem gives
\begin{equation*}
 \|(1-T^l)T^{n_{\rm sl}}\|
 \le\sup_{0\le t\le1}l(1-t)t^{n_{\rm sl}}
 \le l/(n_{\rm sl}+1)\le la_N/(2r).
\end{equation*}
For each normalized trace in~\eqref{eq:E-trace-displacement}, reflection Schwarz at the slab ends and the chessboard inequality bound the remaining factors by the product of their centered roots, at most~$A^n$. One applies this inequality to the complete cyclic product, so no supremum norm of~$X$ is required. The two terms together are bounded by $A^nla_N/r$. An arity-one moment is translation invariant and has change zero.

For a negative move reverse the chosen axis. Coordinate separation is unchanged and the assumed reflected roots have the same bound, so the same slab estimate applies. Replace the $4n$ coordinates one at a time. Each starting configuration is a coordinate hybrid and hence has the required separation. Each move is at most $|x-y|\le r$. Summing proves the second bound. The constants depend only on the roots, arity and slab margin, uniformly in~$N,p$.
\end{proof}
